% report/report.tex
\documentclass[11pt]{article}
\usepackage[a4paper,margin=2.5cm]{geometry}
\usepackage{graphicx}
\usepackage{hyperref}
\usepackage{booktabs}

\title{Adversarial Attacks on RAG-based Recommender Systems: \\ Potential Risks and Countermeasures}
\author{AI4SE Project 12}
\date{\today}

\begin{document}
\maketitle

\section{Introduction}
Motivation for studying data-poisoning attacks against retrieval-augmented
recommender systems, and why countermeasures matter.

\section{Related Work}
Summary of the Poison-RAG paper and how this project's attack design follows it.

\section{System Architecture}
Description of the repository pattern (in-memory vs. file persistence),
the retriever, the generator, and how they compose into the RAG pipeline.
Reference the module diagram: \texttt{src/models.py}, \texttt{src/repository/},
\texttt{src/embeddings.py}, \texttt{src/retriever.py}, \texttt{src/generator.py},
\texttt{src/rag\_pipeline.py}.

\section{Attack Design}
Description of \texttt{PoisonRAGAttacker}: how adversarial items are crafted
and injected, and what makes them effective against the retriever.

\section{Defenses}
Description of \texttt{OutlierFilterDefense}, \texttt{MajorityAgreementDefense},
and \texttt{PromptGuardGenerator} (ICL / Chain-of-Thought prompting), and the
\texttt{MetadataTrustDefense} oracle used as an upper bound.

\section{Experimental Setup}
Models compared: Llama, gpt-oss-20B, Qwen, Mistral. Metrics: attack success
rate, hit rate@k, NDCG@k. Table of configurations run.

\section{Results}
\begin{table}[h]
\centering
\begin{tabular}{lcc}
\toprule
Configuration & Attack Success Rate & Notes \\
\midrule
Baseline (no attack) & & \\
Attacked, no defense & & \\
Attacked + outlier filter & & \\
Attacked + majority agreement & & \\
Attacked + prompt guard (CoT) & & \\
\bottomrule
\end{tabular}
\caption{Attack success rate across configurations and models.}
\end{table}

\section{Discussion}
Which defenses were most effective, trade-offs (utility loss vs. protection),
and differences observed across the four models.

\section{Conclusion}
Summary of findings and directions for future work.

\end{document}
